\chapter{\(3\) temporal sheets, nilpotent calculus, and tritified Cartan connection}
\label{chap:tres-hojas-variacion}

\section{Quadratic Family of TRIT}

For \(\tau\in\{+1,0,-1\}\), consider the real algebra

\[
\mathbb A_\tau=\mathbb R[J_\tau]/(J_\tau^2+\tau),
\qquad
J_\tau^2=-\tau I.
\]

The exponential is computed exactly:

\[
e^{tJ_{+1}}=\cos t+J_{+1}\sin t,
\]

\[
e^{tJ_0}=I+tJ_0,
\]

\[
e^{tJ_{-1}}=\cosh t+J_{-1}\sinh t.
\]

The three sheets have differentiated geometric functions:

\begin{center}
\begin{tabular}{ccl}
\toprule
\(\tau\)&\(J_\tau^2\)&r\'egimen\\
\midrule
\(+1\)&\(-I\)&elliptic: closure, return and memory\\
\(0\)&\(0\)&parabolic: tangency, variation, and present\\
\(-1\)&\(+I\)&hyperbolic: aperture, boost, and propagation\\
\bottomrule
\end{tabular}
\end{center}

Let \(\widehat\tau=P_+-P_-\), with \(P_0=I-P_+-P_-\). Then

\[
\widehat\tau^3=\widehat\tau.
\]

An involution that exchanges the outer sheets and preserves the central sheet
satisfies

\[
\mathsf J\widehat\tau\mathsf J^{-1}=-\widehat\tau.
\]

The threshold remains fixed and the closing and opening orientations are interchanged.

\section{The null sheet as 1st jet}

In \(\mathbb A_0\), write \(\epsilon=J_0\), with \(\epsilon^2=0\). For
every differentiable functional \(F\),

\begin{equation}
\boxed{F(x+\epsilon v)=F(x)+\epsilon\,dF_x(v).}
\label{eq:rm-dual-first-jet}
\end{equation}

The equality is exact because all terms of order at least two
contain \(\epsilon^2\). In particular, for an action
\(S[e,\omega,\Psi]\),

\begin{equation}
S[e+\epsilon\delta e,\omega+\epsilon\delta\omega,
  \Psi+\epsilon\delta\Psi]
=S[e,\omega,\Psi]+\epsilon\,\delta S.
\label{eq:rm-action-first-jet}
\end{equation}

The nilpotent coefficient contains the Euler--Lagrange equations, the
Cartan equation, the symplectic potential, and the boundary terms.
The central sheet is therefore the algebraic residence of the variational
principle.

\begin{theorem}[Nilpotent stationarity]
For a configuration \(z\), the following are equivalent:
\[
dS_z(\delta z)=0\quad\text{for every }\delta z,
\]
and
\[
S(z+\epsilon\delta z)=S(z)
\quad\text{in }\mathbb A_0
\quad\text{for every }\delta z.
\]
\end{theorem}

\begin{proof}
The equivalence is the equality of the coefficients of \(1\) and \(\epsilon\) in
\eqref{eq:rm-dual-first-jet}.
\end{proof}

The global Hamiltonian, the effective action of the central sheet, and the
modular Hamiltonian of the boundary are coordinated objects:

\[
H_{\rm glob},
\qquad
L_{0,\rm eff},
\qquad
K_{\rm HHM}=-\log\rho_{\partial}.
\]

The first generates evolution; the second publishes
local variation; the third preserves informational provenance.

\section{Feshbach--Schur reduction}

Let

\[
\mathcal H=\mathcal H_0\oplus\mathcal H_{\rm ext},
\qquad
H=
\begin{pmatrix}
H_{00}&V\\ V^*&H_{\rm ext}
\end{pmatrix}.
\]

For \(z\) in the resolvent set of \(H_{\rm ext}\), the resolvent
compression is

\begin{equation}
P_0(z-H)^{-1}P_0
=\left[z-H_{00}-\Sigma_0(z)\right]^{-1},
\qquad
\Sigma_0(z)=V(z-H_{\rm ext})^{-1}V^*.
\label{eq:rm-feshbach-schur}
\end{equation}

The operator \(\Sigma_0\) transports the memory of the elliptic and hyperbolic
sheets to the present sheet. Its time transform generates a memory kernel

\[
S_{\rm eff}[\psi_0]
=\int\psi_0^\dagger(i\hbar\partial_t-H_{00})\psi_0\,\mathrm dt
-\iint\psi_0^\dagger(t)K_0(t,t')\psi_0(t')\,\mathrm dt\,\mathrm dt'.
\]

A regular low-frequency expansion,

\[
\Sigma_0(\omega)
=\Sigma_0^{(0)}+\omega\Sigma_0^{(1)}
+\omega^2\Sigma_0^{(2)}+\cdots,
\]

produces an effective local action; a singular spectral structure preserves
nonlocal memory. The present thus receives the response of the global state without
introducing an external time.

\section{Connection universal of Cartan}

Let \(\mathfrak g_\tau\) be the Cartan algebra generated by
\(J_{ab}=-J_{ba}\) and \(P_a^{(\tau)}\), with relations

\begin{align}
 [J_{ab},J_{cd}]&=
 \eta_{bc}J_{ad}-\eta_{ac}J_{bd}
 -\eta_{bd}J_{ac}+\eta_{ad}J_{bc},\\
 [J_{ab},P_c^{(\tau)}]&=
 \eta_{bc}P_a^{(\tau)}-\eta_{ac}P_b^{(\tau)},\\
 [P_a^{(\tau)},P_b^{(\tau)}]&=-\tau J_{ab}.
 \label{eq:rm-cartan-brackets}
\end{align}

For \(\tau=0\) one obtains the algebra of Poincare; the extreme signs yield the two constant-curvature forms. Let \(\omega\) be a Lorentz connection, \(e=e^aP_a^{(\tau)}\) a coframe valued in the translation sector, and \(\ell\) a length scale. We define the tritified connection

\begin{equation}
\mathcal A_\tau^{\rm Cart}
=\frac12\omega^{ab}J_{ab}+\ell^{-1}e^aP_a^{(\tau)}.
\label{eq:rm-trit-cartan-connection}
\end{equation}

Its curvature is

\begin{equation}
\boxed{
\mathcal F_\tau^{\rm Cart}
=\frac12\left(F_\omega^{ab}-\tau\ell^{-2}e^a\wedge e^b\right)J_{ab}
+\ell^{-1}T^aP_a^{(\tau)},
\qquad T^a=D_\omega e^a.}
\label{eq:rm-trit-cartan-curvature}
\end{equation}

\begin{proof}
Upon expanding
\(\mathrm d\mathcal A_\tau^{\rm Cart}
+\mathcal A_\tau^{\rm Cart}\wedge\mathcal A_\tau^{\rm Cart}\), the first
bracket of \cref{eq:rm-cartan-brackets} produces \(F_\omega\), the second
produces \(D_\omega e\), and the third contributes
\(-\tau\ell^{-2}e^a\wedge e^bJ_{ab}/2\).
\end{proof}

In the central sheet,

\[
\mathcal F_0^{\rm Cart}
=\frac12F_\omega^{ab}J_{ab}+\ell^{-1}T^aP_a^{(0)}.
\]

Torsion is the nilpotent coefficient of curvature. This equality
joins the variational interpretation of the sheet
\(0\) to its geometric Cartan
interpretation.

\section{Nonadic corona and torsional energy}

On the edges of the crown of nines, let be

\[
(D_{9,m}f)(a)=f(t(a))-U_a f(s(a)),
\]

where \(U_a\) transports orientation and carry. For a
positive weight \(W_m\),

\begin{equation}
H_m=H_{{\rm clk},m}+D_{9,m}^*W_mD_{9,m},
\qquad
\mathcal E_{{\rm tor},m}(f)
=\|W_m^{1/2}D_{9,m}f\|^2.
\label{eq:rm-nonadic-torsion-energy}
\end{equation}

Unitary holonomy preserves positivity. The refinement squares

\[
D_{9,m+1}J_m=K_mD_{9,m},
\qquad
K_m^*W_{m+1}K_m=W_m
\]

guarantee the compatibility of energy at every depth.

The subsequent closure will construct the screen soldering that transforms this
defect into a Cartan response. The present section fixes the operator
domain and the positivity that this intertwiner must preserve.

\section{Principle of Ambivalence: bidirectional involution of the ratios \(\pi/\varphi^2\) and \(\varphi/\pi^2\)}

We define

\[
\mathcal Q(x,y)=\left(\frac{x}{y^2},\frac{y}{x^2}\right),
\qquad x,y>0.
\]

For

\[
P=xy,
\qquad
O=x/y,
\]

we verify

\[
P(\mathcal Q(x,y))=P^{-1},
\qquad
O(\mathcal Q(x,y))=O^3.
\]

After two steps,

\begin{equation}
\boxed{P(\mathcal Q^2)=P,
\qquad O(\mathcal Q^2)=O^9.}
\label{eq:rm-qmap-nine}
\end{equation}

In coordinates \(u=\log(xy)\), \(v=\log(x/y)\),

\[
(u,v)\longmapsto(-u,3v)\longmapsto(u,9v).
\]

The product returns, and orientation multiplies its memory by nine. For
\((x,y)=(\pi,\varphi)\), the two outputs are

\[
\rho_A=\frac{\pi}{\varphi^2},
\qquad
\rho_E=\frac{\varphi}{\pi^2}.
\]

Satisfy

\[
\rho_E\rho_A^2=\varphi^{-3},
\qquad
M_9=(\rho_E\rho_A^2)^6=\varphi^{-18}.
\]

The relation links areal reading, electronic exponent, nonadic return
and quadratic memory without identifying their algebraic structures with each other.

\begin{proofstatusbox}
The quadratic algebras, the dual calculus, the Schur reduction, the
curvature \eqref{eq:rm-trit-cartan-curvature}, and the identity
\eqref{eq:rm-qmap-nine} are exact. The limit of the torsional energy and its
realization as the full Einstein--Cartan action require the
refinement squares and the physical morphism developed subsequently.
\end{proofstatusbox}

\begin{falsifierbox}
The following refute internal closure: a term other than
\(-\tau\ell^{-2}e\wedge e\) in the curvature; a nilpotent coefficient that is not
\(D_\omega e\); a residue in the Feshbach--Schur identity; loss
of positivity of \(D_9^*WD_9\); or failure of either of the two equations
of \eqref{eq:rm-qmap-nine}.
\end{falsifierbox}
